\begin{abstract}
An implemented sturgeon welfare-monitoring system for recirculating aquaculture pairs a locally
hosted LLM with a deterministic rule engine holding final, escalation-only authority over severity
decisions, because a dissolved-oxygen crash can kill stock within hours and a silent model mistake
is not tolerable. One might expect the architecture's value to come from the rule engine correcting
the model's occasional wrong call. Measured against the system, offline, against synthetically generated sensor
readings and labeled imagery, that is not what happens:
half the evaluated scenarios produce output the model never parses, and the rule engine's
contribution is defaulting safe rather than adjudicating a disagreement -- genuine correction occurs
in only one case in eight. Even where the model's final answer agreed with the rule engine, its
reasoning fabricated a sensor value never present in the input in two of three such cases, invisible
to output-only evaluation. This was tested against a self-documented, insufficient training adapter,
making it a stronger stress test, not a weaker one. The same narrow-interface discipline held across
four further layers: an agentic tool-routing negative result surviving an adversarial re-run, a
retrieval threshold tuned for asymmetric cost, a drift-detection binning choice whose sensitivity
crosses over rather than behaving monotonically, and a vision pipeline whose configured operating
point runs at recall 0.782, not the headline academic-optimum recall of 0.719 usually quoted for
it. The resulting rule is narrow and falsifiable: such a pipeline can be made safe against an
unreliable model provided its validator reads nothing beyond one enumerable field and always
defaults conservatively -- a claim about interface design, not about having made the model
trustworthy.
\end{abstract}
